\section{Compatible choices and the conclusion}
\label{sec:completion}

We choose parameters satisfying all the preceding estimates, then combine
the resulting integral detector and reduced Bott vanishing to prove
Theorem~\ref{thm:main}.  The two indices have different roles:
$i\geq0$ is a height in one infinite
tower; $\nu\geq1$ indexes a sequence of increasingly accurate
constructions.  Degrees are independent of $i$.  No uniform degree
bound in $\nu$ is required.

\subsection{The order of parameters}

First fix the integer $b$ with
$\sqrt{5\cdot2^{1-b}}<3/5$.  This gives the height-independent gap
$\rho=1-.05^2/2$ of Proposition~\ref{prop:sequential-law}.
Retain the explicit compact Bott projection in
\eqref{eq:compact-bott-profile}, its fixed signed class, and
its matrix size.  Fix enlarged support profiles for all neighboring
scales; their ratios depend on $b$, which remains fixed for every
$\nu$.  The rescaling and rotation paths therefore have one common
modulus of continuity.

For each $\nu$, take a target error $\epsilon_\nu\downarrow0$.
Perform the following finite choices in the indicated order.
\begin{enumerate}
\item Choose the constant-extraction length $k_\nu$ from the fixed
gap, so that the term $\rho^{2k_\nu}$ in
\eqref{eq:constant-extraction-error} is sufficiently small.
Choose finite meshes of the two projection paths.  The four endpoint
families, those mesh profiles, and their cutoffs specify a finite
collection of same-height and adjacent-height types.  Choose all
per-factor error tolerances small enough that the telescoping bound
in Lemma~\ref{lem:finite-profile-realization} is below
$\epsilon_\nu$.
\item Choose the finite Fourier truncations and finite matrix-weight
net of Lemma~\ref{lem:product-sampler}.  Its support exponents,
the finite regular-norm test exponent, and the moment constant in
Lemma~\ref{lem:centered-voltage-block} are now fixed.  Also fix a
desired deletion fraction small enough that
$k_\nu\sqrt{\varepsilon_{\mathrm{del}}}$ has the required
error bound.
\item Choose $s_\nu$ as a sufficiently large multiple of $20000$
and $5b$.  Besides the coefficient, probability, and moment
requirements, impose
\[
 \frac{s_\nu}{b}\geq10^{12},\qquad
 \frac b{.50005s_\nu}\leq10^{-6},\qquad
 2^{-.000025s_\nu+O(\log s_\nu)+O_\nu(1)}
       \ll\epsilon_\nu.
\]
The last estimate uses the actual refined alphabet, after the
finite type and weight choices.  Its exponent has a strictly
negative linear term, so it can be imposed along with every
previous finite requirement.  The uniform gate budget and the
height-separation term depending only on $s_\nu,b$ can be made
smaller than the allocated deletion budget at this step.
\item Choose $n_{0,\nu}$, hence
$H_{0,\nu}=(2n_{0,\nu}+1)s_\nu$, sufficiently large.  Impose
$H_{0,\nu}/s_\nu\geq10^{12}$ and all initial-height requirements
in Sections~\ref{sec:labeling}--\ref{sec:assembly}.  In
particular make the remaining separation budget small, make the
sum of the expected pruning fractions smaller than its allocation,
and make the sum of all voltage and norm failure probabilities
less than one.  Increase $H_{0,\nu}$ further to make both the
domain-trimming error and the physical-synthesis error less than
their allocations.  Their bounds have the forms
\[
 C_\nu2^{-345H_{0,\nu}},\qquad
 C_\nu\frac{2^{-344.5H_{0,\nu}}}{1-2^{-344.5b}},
\]
where $C_\nu$ is already independent of height and of this final
increase.  Thus these requirements are compatible.
\end{enumerate}

We explain the choices of random data after these parameters.
Choose the coefficient pools once as in
Lemma~\ref{lem:coefficient-pools}; the same finite pools are reused
on all full blocks.  The sequential law then gives expanding
auxiliary incidences for every table realization in its support.
The deletion estimates in Proposition~\ref{prop:tower-labels}
hold for every such realization.  Proposition~\ref{prop:pruning}
selects a realization with all further layer deletions small.
Only after fixing this old graph apply
Proposition~\ref{prop:sampled-symbols}.  Its one independent
voltage-and-weight law controls both kinds of failure before
their simultaneous avoidance.  There is consequently no change
of measure hidden in the centered moment estimate.

The finite label set is unchanged as heights increase.  It is
refined only by the chosen finite weight net.  The private
completion in Lemma~\ref{lem:private-completion} adds a countable
alphabet, whose coefficients in every operator polynomial are
zero.  The resulting graphical group $\Gamma_\nu$ is countable;
finite generation is obtained from a finite $K$-theory witness
after the quotient argument.

This order also avoids a potential circularity in the norm
sequence quotient.  At stage $\nu$ only finitely many profiles
are required.  Their meshes become finer with $\nu$, and their
uniform modulus of continuity supplies the complete paths in
the quotient.  One does not choose a separate nonuniform
approximation for each height or each path parameter.

\subsection{Proof of the main result}

\begin{proof}[Proof of Theorem~\ref{thm:main}]
Make the sequence of compatible constructions just described.
Proposition~\ref{prop:graphical-realization} embeds $\HH$ in
every $\Gamma_\nu$.
Proposition~\ref{prop:torsion-detector} proves torsion-freeness
and provides an integral class whose evaluation on the ordered
torus $(z,x)$ is $-1$.
Proposition~\ref{prop:reduced-bott-vanishing} proves that, for
all sufficiently large $\nu$, the rank-zero Bott class of this
torus is zero in $K_0(\Cr(\Gamma_\nu))$.

Fix one such $\nu$.  Apply Lemma~\ref{lem:finite-subgroup-witness}
with $H=\HH$ and the physical Bott projections to obtain a finitely
generated subgroup
$\HH\leq \Ginj\leq\Gamma_\nu$ in which that equality already
holds.  The subgroup $\Ginj$ is torsion-free.  Restrict the negative
of the detector class to $B\Ginj$, and denote it by $c$.  Let
$i:\Z^2\hookrightarrow \Ginj$ be the inclusion with ordered
generators $(z,x)$.  The detector evaluation is unchanged by restriction,
so
\[
 \ip{(Bi)^*c}{[T^2]}=1,
 \qquad i_*(\beta)=0\quad\hbox{in }K_0(\Cr(\Ginj)).
\]
It remains to detect the corresponding geometric class and identify its
assembly image.  We use compactly supported $K$-homology throughout. On finite CW
complexes, the natural comparison of geometric cycles with analytic
$K$-homology sends $(M,E,f)$ to $f_*[D_E]$
\cite[Theorems~6.1--6.2]{BHS07}; passage to finite-subcomplex colimits
gives the convention used for $B\Ginj$. In particular the compact torus
cycle used here defines a class without any compactness assumption on
$B\Ginj$.

Let $L$ be the line bundle on $B\Ginj$ classified by $c$. The naturality
of the index pairing follows at a finite stage from pullback of vector
bundles and associativity of the Kasparov product:
\[
 [L]\otimes f_*[D]=[f^*L]\otimes[D],\qquad f=Bi.
\]
The same identity holds for the virtual bundle $[L]-[1]$. Applying the
twisted spin Dirac index theorem \cite{AS63}, in the dimension-two form
of \cite[Theorem~1]{BvE16}, gives
\begin{align*}
 \bigl\langle [L]-[1],(Bi)_*[D_{\mathbb T^2}]\bigr\rangle
 &=\operatorname{index}(D_{\mathbb T^2}\otimes (Bi)^*L)
       -\operatorname{index}(D_{\mathbb T^2})\\
 &=\bigl\langle (Bi)^*c,[\mathbb T^2]\bigr\rangle=1.
\end{align*}
The torus tangent bundle is trivial, its $\widehat A$ class is one,
and its untwisted complex Dirac index is zero. Thus the pushed-forward
$K$-homology class has infinite order and is nonzero after tensoring with
$\Q$.

Under Fourier transform $\Cr(\Z^2)\cong C(\mathbb T^2)$, the assembly
image of the torus spin Dirac class is the rank-zero Bott generator, up to
orientation. This is the full-torus case of Emerson--Hudson's assembly
formula for a $K$-oriented subtorus \cite[Theorem~1.3]{EH20}: the dual
subtorus is a point, whose shriek class is the compact Bott class.

For completeness we verify the required subgroup naturality at this
particular cycle. Put $A=\Cr(\Z^2)$, $B=\Cr(\Ginj)$, and write
$j:A\to B$ for the isometric inclusion. If $P\to T^2$ is the principal
$\Z^2$-bundle classified by the torus map, its Mi\v{s}\v{c}enko bundle
is $E_A=P\times_{\Z^2}A$. The bundle obtained after the subgroup map is
canonically
\[
 (P\times_{\Z^2}\Ginj)\times_{\Ginj}B
 \cong P\times_{\Z^2}B
 \cong E_A\otimes_j B.
\]
The last identification sends $([p,a],b)$ to $[p,j(a)b]$; the balancing
relations and fiber inner products agree because $j$ sends each subgroup
unitary to the same group unitary. Extension of coefficients therefore
commutes with the index Kasparov product. In explicit notation,
\[
 \bigl([E_A]\otimes_{C(T^2)\otimes A}([D_{T^2}]\boxtimes1_A)\bigr)
      \otimes_A[j]
 =[E_A\otimes_j B]\otimes_{C(T^2)\otimes B}
      ([D_{T^2}]\boxtimes1_B).
\]
These indices are the assembly classes by the Mi\v{s}\v{c}enko comparison
\cite[Corollary~1.4 and diagram~(1.6)]{KP22}; see also
\cite[Proposition~3.7 and Theorem~4.1]{Land15}.
The comparison applies to the compact torus cycle, and its
finite-subcomplex colimit handles the target $B\Ginj$.
Consequently
\[
 \mu_{\Ginj}^r\bigl((Bi)_*[D_{T^2}]\bigr)
   =j_*\mu_{\Z^2}^r([D_{T^2}])=0
\]
by Proposition~\ref{prop:reduced-bott-vanishing} and
Lemma~\ref{lem:finite-subgroup-witness}. This is integral vanishing,
whereas the index pairing above proves infinite order of the source class.
\end{proof}
